\section{Conclusion}
\label{sec:conclusion}

This paper presented a systematic measurement study comparing SSHv2, SSH3, and Mosh as transport substrates for interactive AI agent workloads. Using a controlled testbed with tc/netem emulation across three network scenarios, we evaluated session establishment, command completion latency, and output completeness.

Our key findings are as follows. First, SSH3 achieves 2.2--3.4$\times$ faster session establishment than SSHv2 and 3.2--3.7$\times$ faster than Mosh, owing to QUIC's 1-RTT handshake (RQ1). Second, for lightweight commands ($<$1\,KB output), SSH3 provides the lowest latency under impairment (1.5$\times$ faster than SSHv2 in the High scenario). However, for heavy-output commands, SSH3's userspace QUIC is 1.7--2.2$\times$ slower than SSHv2's kernel TCP, which benefits from decades of optimization (CUBIC, pacing, TSO/GRO offload) (RQ2, RQ3). Third, Mosh achieves the lowest apparent latency for all commands but delivers only 2--6\% of expected output for heavy commands and 89--92\% for lightweight commands, due to SSP's frame-skipping design. For AI agents that must parse complete output, this renders Mosh unsuitable for most workloads (RQ4).

Based on these results, we offer the following practical guidelines for AI agent deployments: (1)~when session churn is frequent (e.g., per-task connections), SSH3 provides the lowest setup overhead; (2)~for commands producing small output (status checks, simple queries), SSH3 offers the best latency; (3)~for commands producing large output (log retrieval, directory listings), SSHv2 remains the most efficient choice; (4)~Mosh should only be used for interactive editing tasks where output completeness is not required.

\textbf{Limitations.} Our measurements use a single-hop controlled testbed with emulated impairments. Real-world deployments may encounter middlebox interference (particularly for UDP-based protocols), multi-hop routing variability, and server load effects not captured here. SSH3 is still experimental software (v0.1.7) and its QUIC performance may improve as the implementation matures.

\textbf{Future work.} We plan to extend this study to wide-area networks, evaluate the impact of 0-RTT resumption in SSH3, and integrate our benchmark framework with actual AI agent systems to measure end-to-end task completion time.
